% =====================================================================
%  Apresentacao - Despacho Hidrotermico com Unit Commitment
%  Otimizacao Combinatoria
%
%  Compilar:  pdflatex apresentacao.tex   (rodar 2x para o sumario)
%
%  PLACEHOLDERS: procure por "% PREENCHER" -- sao as tabelas de
%  resultados que voce completa depois de rodar os experimentos.
% =====================================================================
\documentclass[10pt,aspectratio=169]{beamer}

\usepackage[utf8]{inputenc}
\usepackage[T1]{fontenc}
\usepackage[brazil]{babel}
\usepackage{amsmath,amssymb}
\usepackage{booktabs}
\usepackage{graphicx}
\usepackage{array}

\usetheme{Boadilla}
\usecolortheme{seahorse}
\setbeamertemplate{navigation symbols}{}
\setbeamertemplate{caption}[numbered]
\setbeamerfont{block title}{size=\small}

% ---- Notacao ----
\newcommand{\Nb}{\mathcal{N}}          % barras
\newcommand{\Lb}{\mathcal{L}}          % linhas
\newcommand{\Hb}{\mathcal{H}}          % UHEs
\newcommand{\Gb}{\mathcal{G}}          % UTEs
\newcommand{\Tb}{\mathcal{T}}          % etapas
\newcommand{\pH}{p^{H}}                % geracao hidraulica
\newcommand{\pT}{p^{T}}                % geracao termica
\newcommand{\uH}{u^{H}}                % estado UHE
\newcommand{\uT}{u^{T}}                % estado UTE

% ---- Displays compactos (slides densos de formulacao) ----
% Use DEPOIS de \small: mudanca de corpo de fonte reseta estes comprimentos.
\newcommand{\displaycompacto}{%
  \setlength{\abovedisplayskip}{1pt}%
  \setlength{\belowdisplayskip}{1pt}%
  \setlength{\abovedisplayshortskip}{1pt}%
  \setlength{\belowdisplayshortskip}{1pt}%
}

\title[Despacho Hidrotérmico]{Despacho Hidrotérmico com \emph{Unit Commitment}}
\subtitle{Branch \& Bound, Planos de Corte e Relaxação Lagrangeana}
\author{Guilherme Bassous}
\institute{Otimização Combinatória I}
\date{08 de setembro de 2026}

\begin{document}

% ===================== 1. CAPA =====================
\begin{frame}[plain]
  \titlepage
\end{frame}

% ===================== 2. INTRODUCAO =====================
\begin{frame}{O problema}

  Programar a operação de um sistema hidrotérmico ao longo de um horizonte
  discreto, minimizando custo de combustível, partidas e déficit, respeitando
  a rede elétrica e a dinâmica dos reservatórios.

  \vspace{0.4em}
  \begin{columns}[T]
    \begin{column}{0.48\textwidth}
      \begin{block}{Decisões}
        \footnotesize
        \begin{itemize}\itemsep2pt
          \item \textbf{Contínuas:} geração hidráulica e térmica, turbinamento,
                vertimento, defluência, armazenamento, fluxos, ângulos, déficit
          \item \textbf{Binárias:} estado das usinas (ligada/desligada),
                partida e desligamento das térmicas
        \end{itemize}
      \end{block}
    \end{column}
    \begin{column}{0.48\textwidth}
      \begin{block}{Restrições}
        \footnotesize
        \begin{itemize}\itemsep2pt
          \item Atendimento à demanda (1\textsuperscript{a} lei de Kirchhoff)
          \item Fluxo DC e limites de transmissão (2\textsuperscript{a} lei)
          \item Balanço hídrico e cascata entre reservatórios
          \item Capacidade e geração mínima (acopladas às binárias)
          \item \alert{Rampas} de geração e de defluência
          \item \alert{Tempo mínimo} ligado / desligado
        \end{itemize}
      \end{block}
    \end{column}
  \end{columns}

  \vspace{0.3em}
  \begin{block}{Escala da instância --- 2 semanas em base horária}
    \centering\scriptsize
    \begin{tabular}{lr@{\qquad}lr@{\qquad}lr}
      \toprule
      Barras ($N$) & 14  & Variáveis contínuas & $24\,528$ & Restrições         & $34\,573$ \\
      Linhas ($L$) & 20  & Variáveis binárias  & $6\,384$  & \quad rampas       & $6\,030$  \\
      UHEs ($H$)   & 4   & \quad \emph{total}  & $30\,912$ & \quad tempo mínimo & $3\,339$  \\
      UTEs ($G$)   & 5   & Etapas ($T$)        & 336       & \quad rede (KCL/KVL) & $11\,424$ \\
      \bottomrule
    \end{tabular}
  \end{block}
\end{frame}

% ===================== 3. SISTEMA-TESTE =====================
\begin{frame}{Sistema-teste: IEEE 14 barras}
  \begin{columns}[T]
    \begin{column}{0.52\textwidth}
      \footnotesize
      \textbf{Rede.} 14 barras, 20 linhas, reatâncias em p.u.\ na base de
      100~MVA. Susceptância $B_{\ell\ell} = S_{base}/x_\ell$, de modo que o
      fluxo $f = B A^\top \theta$ já sai em MW e a 2\textsuperscript{a} lei
      distribui os fluxos pelas malhas.

      \vspace{0.6em}
      \textbf{Hidráulico.} Cascata
      $\text{UHE}_1 \to \text{UHE}_2 \to \text{UHE}_3 \to \text{UHE}_4$,
      com defluência mínima e limite de deplecionamento.

      \vspace{0.6em}
      \textbf{Térmico.} Cinco máquinas, da base barata
      ($42$~R\$/MWh) ao \emph{peaker} caro ($450$~R\$/MWh), com custo de
      partida, rampa e tempos mínimos distintos.

      \vspace{0.6em}
      \textbf{Déficit.} Penalizado a $3\,000$~R\$/MWh em cada barra, o que
      garante viabilidade do despacho.
    \end{column}
    \begin{column}{0.46\textwidth}
      \centering
      \IfFileExists{figuras/1_geracao_demanda.png}{%
        \includegraphics[width=\linewidth]{figuras/1_geracao_demanda.png}\\
        {\tiny Geração por fonte \emph{vs.} demanda total}%
      }{%
        \fbox{\parbox[c][3.2cm][c]{0.9\linewidth}{\centering\footnotesize
          \emph{figuras/1\_geracao\_demanda.png}}}%
      }

      \vspace{0.8em}
      \scriptsize
      \begin{tabular}{lrrr}
        \toprule
        UTE & $\bar p$ & CVU & $t^{on}/t^{off}$ \\
        \midrule
        1 & 110 & 42  & 6 / 5 \\
        2 &  85 & 85  & 4 / 3 \\
        3 &  70 & 130 & 3 / 2 \\
        4 &  60 & 210 & 2 / 2 \\
        5 &  45 & 450 & 2 / 2 \\
        \bottomrule
      \end{tabular}
    \end{column}
  \end{columns}
\end{frame}

% ===================== 4. VARIAVEIS =====================
\begin{frame}{Formulação: variáveis de decisão}
  \small
  Índices: barra $n \in \Nb$, linha $\ell \in \Lb$, UHE $u \in \Hb$,
  UTE $g \in \Gb$, etapa $t \in \Tb$.

  \vspace{0.5em}
  \centering
  \begin{tabular}{llll}
    \toprule
    Símbolo & Descrição & Tipo & Código \\
    \midrule
    $\pH_{u,t}$    & geração da UHE $u$                & $\mathbb{R}_+$    & \texttt{g\_uhe} \\
    $q_{u,t}$      & turbinamento                      & $\mathbb{R}_+$    & \texttt{turb\_uhe} \\
    $s_{u,t}$      & vertimento                        & $\mathbb{R}_+$    & \texttt{vert\_uhe} \\
    $d_{u,t}$      & defluência                        & $\mathbb{R}_+$    & \texttt{defl\_uhe} \\
    $v_{u,t}$      & armazenamento                     & $\mathbb{R}_+$    & \texttt{uhe\_arm} \\
    $\uH_{u,t}$    & estado da UHE                     & $\{0,1\}$         & \texttt{uch} \\
    \midrule
    $\pT_{g,t}$    & geração da UTE $g$                & $\mathbb{R}_+$    & \texttt{g\_ute} \\
    $\uT_{g,t}$    & estado da UTE                     & $\{0,1\}$         & \texttt{uct} \\
    $y_{g,t}$      & partida                           & $\{0,1\}$         & \texttt{part\_ute} \\
    $z_{g,t}$      & desligamento                      & $\{0,1\}$         & \texttt{desl\_ute} \\
    \midrule
    $f_{\ell,t}$   & fluxo na linha $\ell$             & $\mathbb{R}$      & \texttt{fluxo} \\
    $\theta_{n,t}$ & ângulo de fase                    & $\mathbb{R}$      & \texttt{theta} \\
    $r_{n,t}$      & déficit na barra $n$              & $\mathbb{R}_+$    & \texttt{deficit} \\
    \bottomrule
  \end{tabular}
\end{frame}

% ===================== 5. FO + REDE =====================
\begin{frame}{Formulação: objetivo e rede elétrica}
  \begin{block}{Função objetivo}
    \[
      \min \;\; \underbrace{\sum_{g \in \Gb}\sum_{t \in \Tb} c_g \, \pT_{g,t}}_{\text{combustível}}
      \;+\; \underbrace{\sum_{g \in \Gb}\sum_{t \in \Tb} c^{\text{par}}_g \, y_{g,t}}_{\text{partidas}}
      \;+\; \underbrace{c^{\text{def}} \sum_{n \in \Nb}\sum_{t \in \Tb} r_{n,t}}_{\text{déficit}}
    \]
  \end{block}

  \begin{block}{1\textsuperscript{a} lei de Kirchhoff --- atendimento à demanda}
    \[
      M^H \pH_t \;+\; M^T \pT_t \;+\; r_t \;-\; A f_t \;=\; D_t,
      \qquad \forall t \in \Tb
    \]
    \footnotesize
    $A \in \{-1,0,1\}^{N \times L}$ é a matriz de incidência barra--linha
    ($+1$ na origem, $-1$ no destino); $M^H, M^T$ localizam as usinas nas barras.
  \end{block}

  \begin{block}{2\textsuperscript{a} lei de Kirchhoff --- fluxo DC}
    \[
      f_t = B A^\top \theta_t, \qquad
      \theta_{\text{ref},t} = 0, \qquad
      -\bar f_\ell \le f_{\ell,t} \le \bar f_\ell
    \]
  \end{block}
\end{frame}

% ===================== 6. HIDRAULICO =====================
\begin{frame}{Formulação: subsistema hidráulico}
  \small
  \displaycompacto
  \begin{block}{Balanço hídrico e produção}
    \begin{align*}
      q_{u,t} + s_{u,t} &= d_{u,t}                 & \forall u,t \\
      \pH_{u,t}         &= \rho_u \, q_{u,t}       & \forall u,t
    \end{align*}
  \end{block}

  \begin{block}{Dinâmica dos reservatórios (cascata)}
    \begin{align*}
      v_{u,1} &= \tfrac{1}{2}\bar v_u + a_{u,1} - q_{u,1} - s_{u,1}
                 + \textstyle\sum_{m \in \text{mont}(u)} d_{m,1} & \forall u \\
      v_{u,t} &= v_{u,t-1} + a_{u,t} - q_{u,t} - s_{u,t}
                 + \textstyle\sum_{m \in \text{mont}(u)} d_{m,t} & \forall u,\; t \ge 2 \\
      v_{u,T} &\ge \tfrac{1}{2}\bar v_u                          & \forall u
    \end{align*}
    \footnotesize
    $a_{u,t}$ é a vazão incremental; o armazenamento final evita esvaziar
    os reservatórios ao fim do horizonte.
  \end{block}

  \begin{block}{Comprometimento e deplecionamento}
    \begin{align*}
      \underline{p}^H_u \, \uH_{u,t} \;\le\; \pH_{u,t} &\;\le\; \bar p^H_u \, \uH_{u,t}
        & \forall u,t \\
      d_{u,t} \;\ge\; \underline{d}_u, \qquad
      \bigl| d_{u,t} - d_{u,t-1} \bigr| &\;\le\; \Delta^H_u
        & \forall u,\; t \ge 2
    \end{align*}
  \end{block}
\end{frame}

% ===================== 7. TERMICO =====================
\begin{frame}{Formulação: subsistema térmico}
  \begin{block}{Capacidade acoplada ao estado}
    \[
      \underline{p}^T_g \, \uT_{g,t} \;\le\; \pT_{g,t} \;\le\; \bar p^T_g \, \uT_{g,t},
      \qquad \forall g \in \Gb,\; t \in \Tb
    \]
    \footnotesize
    Máquina desligada gera zero; ligada, opera acima do mínimo técnico.
    É a relação \emph{variable upper/lower bound} que o solver explora em
    \emph{implied bound cuts}.
  \end{block}

  \begin{block}{Rampas de tomada e alívio de carga}
    \[
      -R_g \;\le\; \pT_{g,t} - \pT_{g,t-1} \;\le\; R_g,
      \qquad \forall g \in \Gb,\; t \ge 2
    \]
  \end{block}

  \vspace{0.3em}
  \begin{alertblock}{Observação sobre os dados}
    \footnotesize
    Para a UTE\,1 vale $\underline{p}^T_1 = 45 > R_1 = 22$: partir a máquina
    exigiria um salto de $0$ para $45$~MW, que viola a rampa. Na prática
    a UTE\,1 fica presa ao estado inicial --- modelos reais usam um limite
    de rampa \emph{de partida} separado.
  \end{alertblock}
\end{frame}

% ===================== 8. COMPROMETIMENTO =====================
\begin{frame}{Formulação: \emph{unit commitment} e tempos mínimos}
  \small
  \displaycompacto
  \begin{block}{Lógica de transição de estado}
    \begin{align*}
      \uT_{g,1} - \uT_{g,0} &= y_{g,1} - z_{g,1}   & \forall g \\
      \uT_{g,t} - \uT_{g,t-1} &= y_{g,t} - z_{g,t} & \forall g,\; t \ge 2 \\
      y_{g,t} + z_{g,t} &\le 1                     & \forall g,t
    \end{align*}
    \footnotesize
    $\uT_{g,0}$ é o estado herdado da etapa anterior ao horizonte.
  \end{block}

  \begin{block}{Tempo mínimo ligado / desligado (janelas contíguas)}
    \begin{align*}
      \sum_{s = t - t^{on}_g + 1}^{t} y_{g,s} &\;\le\; \uT_{g,t}
        & \forall g,\; t \ge t^{on}_g \\[2pt]
      \sum_{s = t - t^{off}_g + 1}^{t} z_{g,s} &\;\le\; 1 - \uT_{g,t}
        & \forall g,\; t \ge t^{off}_g
    \end{align*}
    \footnotesize
    Leitura: se houve partida nas últimas $t^{on}_g$ etapas, a máquina
    \emph{tem} de estar ligada agora --- e simetricamente para o desligamento.
    São $3\,339$ restrições, o núcleo combinatório do problema.
  \end{block}
\end{frame}

% ===================== 9. B&B CONFIGURACAO =====================
\begin{frame}{Branch \& Bound: configuração experimental}
  \small
  Para isolar o efeito de cada componente, o solver foi configurado com os
  recursos automáticos desligados:

  \vspace{0.6em}
  \centering
  \begin{tabular}{llp{5.6cm}}
    \toprule
    Parâmetro & Valor & Motivo \\
    \midrule
    \texttt{Threads}       & 1 & determinismo e comparação justa de tempos \\
    \texttt{Presolve}      & 0 & preservar a formulação original \\
    \texttt{Heuristics}    & 0 & isolar a busca do efeito das heurísticas \\
    \texttt{Cuts}          & 0/1 & desligado no \emph{baseline} \\
    \texttt{InfProofCuts}  & 0 & desliga os cortes \emph{learned} (análise de conflito) \\
    \texttt{PreCrush}      & 1 & necessário para injetar cortes via \texttt{cbCut} \\
    \bottomrule
  \end{tabular}

  \vspace{1em}
  \raggedright
  \begin{block}{Métricas coletadas}
    \footnotesize
    Nós explorados, tempo de resolução, limite inferior na raiz, melhor
    solução inteira encontrada e \emph{gap} final de otimalidade.
  \end{block}
\end{frame}

% ===================== 10. B&B RESULTADOS =====================
\begin{frame}{Branch \& Bound: resultados}
  \centering
  \begin{table}
    \small
    \begin{tabular}{lrrrr}
      \toprule
      Configuração & Nós & Tempo (s) & Melhor sol. (R\$) & Gap alvo \\
      \midrule
      Cortes do solver (padrão)    & $31$       & $11$     & $9\,632\,934{,}7$ & $10^{-4}$ \\
      \emph{Baseline} (sem cortes) & $174\,884$ & $5\,460$ & $9\,624\,105{,}8$ & $10^{-3}$ \\
      \bottomrule
    \end{tabular}
    \caption{Desempenho do \emph{branch \& bound} sob diferentes configurações
      (\texttt{Threads}~1, sem \emph{presolve} e sem heurísticas nos dois casos).}
  \end{table}

  \vspace{0.5em}
  \begin{block}{Comentários}
    \footnotesize
    Com os cortes automáticos ligados o solver separa $658$ cortes
    ($341$ \emph{implied bound}, $266$ \emph{learned}, $34$ Gomory, $12$
    \emph{flow cover}, $5$ MIR) e fecha o gap de $10^{-4}$ em \alert{$31$ nós}.
    Desligando os cortes a árvore explode: mesmo relaxando o alvo para
    $10^{-3}$ são \alert{$174\,884$ nós} e $91$~min --- cerca de $500\times$ mais
    tempo. O limite da relaxação linear é fraco porque as binárias de tempo
    mínimo e partida só se acoplam pelas restrições de rampa; sem cortes o
    \emph{bound} sobe devagar e o gap passa a cair muito lentamente.
  \end{block}
\end{frame}

% ===================== 11. CORTES: MIR AGREGADO =====================
\begin{frame}{Planos de corte I: MIR agregado}
  \small
  Separação de cortes \emph{Mixed Integer Rounding} a partir de combinações
  de linhas do modelo, implementada em \emph{callback} (\texttt{MIPNODE}).

  \vspace{0.4em}
  \begin{block}{Procedimento, por linha-base}
    \footnotesize
    \begin{enumerate}\itemsep2pt
      \item \textbf{Agregar} --- combinar até $k$ linhas para eliminar variáveis
            contínuas estritamente entre seus limites (proxy de variável básica).
      \item \textbf{Transladar} --- levar cada variável ao limite mais próximo
            do ponto relaxado, de modo que todas fiquem $\ge 0$.
      \item \textbf{Aplicar MIR} --- sobre $\sum_j a_j x_j \le b$, com
            $f_0 = b - \lfloor b \rfloor$:
            \[
              \sum_{j \in I} \Bigl(\lfloor a_j \rfloor
                + \tfrac{(f_j - f_0)^+}{1 - f_0}\Bigr) x_j
              \;+\; \sum_{j \in C} \tfrac{\min(a_j,0)}{1-f_0}\, x_j
              \;\le\; \lfloor b \rfloor
            \]
      \item \textbf{Destransladar} e testar violação; repetir para várias
            escalas $\lambda$ da linha, guardando o corte mais violado.
    \end{enumerate}
  \end{block}

  \begin{block}{Filtros de qualidade}
    \footnotesize
    Violação \alert{relativa} ($\text{viol}/\max(1,|b|)$), descarte de cortes
    mal condicionados ($\max|a_j| / \min|a_j| > 10^7$) e \alert{deduplicação}
    por assinatura do corte normalizado.
  \end{block}
\end{frame}

% ===================== 12. CORTES: AGREGACAO 1a LEI =====================
\begin{frame}{Planos de corte II: agregação da 1\textsuperscript{a} lei}
  \small
  Somando a 1\textsuperscript{a} lei sobre \emph{todas} as barras, o termo de
  rede desaparece: cada linha entra com $+1$ na origem e $-1$ no destino, logo
  $\mathbf{1}^\top A = 0$. Sobra o balanço de energia do sistema:
  \[
    \sum_{u \in \Hb} \pH_{u,t} \;+\; \sum_{g \in \Gb} \pT_{g,t} \;+\; S_t \;=\; D_t,
    \qquad S_t = \sum_{n \in \Nb} r_{n,t}, \quad D_t = \sum_{n \in \Nb} D_{n,t}
  \]

  Com $\pH_{u,t} \le \bar p^H_u \uH_{u,t}$ e $\pT_{g,t} \le \bar p^T_g \uT_{g,t}$
  obtém-se uma \alert{mochila} nas binárias de \emph{unit commitment}:
  \[
    \sum_j P_j z_j + S_t \;\ge\; D_t,
    \qquad z_j \in \{0,1\}, \quad S_t \ge 0
  \]

  \begin{block}{Fortalecimento por MIR}
    \footnotesize
    A desigualdade acima já é satisfeita pela relaxação contínua. Complementando
    $z_j = 1 - \bar z_j$, passando a ``$\le$'' e aplicando MIR com
    $b = \lambda(P - D_t)$, $f_0 = b - \lfloor b \rfloor$:
    \[
      \sum_j \bar a_j z_j + \frac{\lambda}{1-f_0}\, S_t
      \;\ge\; \sum_j \bar a_j - \lfloor b \rfloor
    \]
    Para $\lambda = 1$ e $P_j$ inteiro isso vira
    $\;\sum_j P_j z_j + S_t/(1-f_0) \ge \lceil D_t \rceil$: \textbf{arredonda a
    demanda para cima e encarece o déficit}. O corte é separado por etapa, no nó raiz.
  \end{block}
\end{frame}

% ===================== 13. CORTES: RESULTADOS =====================
\begin{frame}{Planos de corte: resultados}
  \centering
  \begin{table}
    \small
    \begin{tabular}{lrrrr}
      \toprule
      Estratégia & Nós & $\Delta$ nós (\%) & Tempo (s) & Melhor sol. (R\$) \\
      \midrule
      Sem cortes (\emph{baseline})                    & $174\,884$ & --        & $5\,460$ & $9\,624\,105{,}8$ \\
      MIR agregado (agrupamento geral)                & $130\,323$ & $-25{,}5$ & $4\,333$ & $9\,624\,057{,}7$ \\
      Agregação 1\textsuperscript{a} lei (por etapa)  & $179\,503$ & $+2{,}6$  & $5\,005$ & $9\,624\,209{,}5$ \\
      \bottomrule
    \end{tabular}
    \caption{Efeito dos planos de corte sobre a árvore (gap alvo $10^{-3}$ em todos).}
  \end{table}

  \vspace{0.3em}
  \raggedright
  \begin{block}{Comentários}
    \footnotesize
    O MIR agregado sobre linhas escolhidas heuristicamente é o mais efetivo:
    corta \alert{$25{,}5\%$ dos nós} e $20{,}6\%$ do tempo em relação ao
    \emph{baseline}. A agregação da 1\textsuperscript{a} lei, apesar de ser um
    corte válido derivado à mão, quase não move a árvore ($+2{,}6\%$ de nós,
    $-8{,}3\%$ de tempo): arredondar a demanda para cima só morde quando
    $D_t$ tem parte fracionária relevante, e a folga de déficit $S_t$ absorve
    boa parte da violação. As três configurações terminam em soluções
    equivalentes dentro da tolerância de $10^{-3}$ (dispersão de $0{,}0016\%$).
  \end{block}
\end{frame}

% ===================== 14. LAGRANGEANA: FORMULACAO =====================
\begin{frame}{Relaxação lagrangeana: formulação}
  \small
  Dualizam-se os blocos que acoplam o problema no tempo dentro de cada máquina:
  \alert{rampas} e \alert{tempo mínimo} ($9\,369$ restrições, todas ``$\le$'',
  logo $\lambda \ge 0$).

  \[
    z^* = \min_{x} \; \bigl\{\, c(x) \;:\; x \in X,\;\; A_R x \le b_R,\;\; A_M x \le b_M \,\bigr\}
  \]
  \[
    L(\lambda) \;=\; \min_{x \in X} \;
      c(x) + \lambda_R^\top (A_R x - b_R) + \lambda_M^\top (A_M x - b_M)
  \]

  \begin{columns}[T]
    \begin{column}{0.5\textwidth}
      \begin{block}{Fica em $X$}
        \footnotesize
        1\textsuperscript{a} e 2\textsuperscript{a} leis, balanço hídrico e
        cascata, capacidade e geração mínima, \emph{lógica de transição}
        $\uT_t - \uT_{t-1} = y_t - z_t$.
      \end{block}
    \end{column}
    \begin{column}{0.5\textwidth}
      \begin{block}{Propriedades}
        \footnotesize
        \begin{itemize}\itemsep1pt
          \item $L(\lambda) \le z^*$ para todo $\lambda \ge 0$
          \item $L$ é côncava e linear por partes
          \item $s = A x^*(\lambda) - b$ é subgradiente de $L$
        \end{itemize}
      \end{block}
    \end{column}
  \end{columns}

  \vspace{0.4em}
  \begin{alertblock}{Dual lagrangeano}
    \centering
    $z_{LD} = \max_{\lambda \ge 0} L(\lambda)$, \qquad com \qquad
    $z_{LP} \;\le\; z_{LD} \;\le\; z^*$
    \\[2pt]
    {\footnotesize $z_{LD} > z_{LP}$ porque $X$ \emph{não} tem a propriedade
    da integralidade.}
  \end{alertblock}
\end{frame}

% ===================== 15. LAGRANGEANA: SUBGRADIENTE E RESULTADOS =====================
\begin{frame}{Relaxação lagrangeana: subgradiente e resultados}
  \begin{columns}[T]
    \begin{column}{0.42\textwidth}
      \small
      \begin{block}{Método do subgradiente}
        \footnotesize
        $\lambda^0 = 0$; a cada iteração:
        \begin{enumerate}\itemsep1pt
          \item resolver o subproblema $\Rightarrow x^k,\, L^k$
          \item $s^k = A x^k - b$
          \item $\alpha_k = \rho_k \dfrac{UB - L^k}{\|s^k\|^2}$ \hfill (Polyak)
          \item $\lambda^{k+1} = \max\bigl(0,\; \lambda^k + \alpha_k s^k\bigr)$
          \item $\rho_k \leftarrow \rho_k/2$ após $p$ iterações sem melhora
        \end{enumerate}
      \end{block}
    \end{column}
    \begin{column}{0.56\textwidth}
      \begin{table}
        \scriptsize
        \begin{tabular}{lrr}
          \toprule
          Limite & Valor (R\$) & Gap p/ $UB$ (\%) \\
          \midrule
          $z_{LP}$ \; relaxação linear   & $9\,594\,580{,}10$ & $0{,}41$ \\
          $z_{LD}$ \; dual lagrangeano   & $9\,618\,970{,}36$ & \alert{$0{,}15$} \\
          $UB$ \;\;\;\, incumbente MIP   & $9\,633\,634{,}50$ & --  \\
          \midrule
          \multicolumn{2}{l}{Ganho $z_{LD} - z_{LP}$}        & $24\,390{,}26$ \\
          \multicolumn{2}{l}{Fração do gap fechada (\%)}     & $62{,}5$ \\
          \multicolumn{2}{l}{Iterações (melhor em $k=0$)}    & $40$ \\
          \bottomrule
        \end{tabular}
        \caption{\scriptsize Qualidade do limite lagrangeano
          ($40$ iterações, subproblema com \emph{gap} $10^{-3}$ e $60$~s).}
      \end{table}
    \end{column}
   \end{columns}
   \begin{alertblock}{Resultado}
     \centering
     $z_{LD} > z_{LP}$ por $24\,390$ (R\$), isto é, \alert{$62{,}5\%$} da
     distância entre a relaxação linear e o incumbente --- confirma que $X$
     não tem a propriedade da integralidade.
   \end{alertblock}

\end{frame}

% ===================== 16. LAGRANGEANA: CONVERGENCIA =====================
\begin{frame}{Relaxação lagrangeana: convergência do subgradiente}
  \centering
  \includegraphics[width=0.94\linewidth]{figuras/5_lagrange_convergencia.png}

  \vspace{-0.1em}
  \raggedright
  \begin{block}{Leitura}
    \footnotesize
    \begin{itemize}\itemsep1pt
      \item Melhor limite: \alert{$L(0) = 9\,618\,970{,}36$}. Já em
            $\lambda = 0$ --- só removendo rampas e tempo mínimo --- o
            subproblema \emph{inteiro} rende $0{,}25\%$ acima de $z_{LP}$, e
            nenhuma das $40$ iterações seguintes supera esse valor.
    \end{itemize}
  \end{block}
\end{frame}

\end{document}
